% HD-R19. GPT-6 Astra Pro /root; inherited effort not exposed.
% 21 November 2026. Actual multiplicity-four source, reflection rank6--8.
\section{HD-R19: the small real factors in the multiplicity-four source}
\label{sec:rank-four-small-real-tensor}

The actual multiplicity-four alternative in HD-I23 has no
still-missing sixfold source with constituent rank six, seven
or eight. Here the coefficient rank is \(m=4r\), the original
group has no nonidentity element of moving rank below four,
and the original sphere coefficient system has parabolic
dimension six and period rank nine. The imprimitive
reflection constituents are already supplied through the
actual rank-four quotient blocks in HD-I23. The primitive
constituents of ranks six through eight are the rational
real models \(A_r,E_6,E_7,E_8\), with the applicable rank.
We retain their exact actual tensor maps and sign kernel.

\paragraph{1. The finite tensor preimage and the original tuple.}
HD-I23 constructs the original-field isomorphism
\[
 C:V_K\xrightarrow{\sim}W_K\otimes_K M_K,
             \qquad \dim W_K=r,\quad\dim M_K=4.          \tag{R19.1}
\]
It uses a rank-one reflection projector \(P\) in the
actual coefficient algebra \(\mathcal A\):
\(\mathcal A P\otimes P V\to V\),
\(aP\otimes u\mapsto au\), with inverse
\(v\mapsto\sum_j a_jP\otimes Pb_jv\) for
\(\sum_j a_jPb_j=1\). The rational real reflection-model
intertwiner is then an actual invertible solution of the
finite \(K\)-linear intertwiner equations. The normalizer
calculation there gives a finite real orthogonal group
\(R^{\pm}\) on \(W\), containing the reflection model
and \(-I\), such that every original element is an
actual tensor \(A\otimes B\) with
\(A\in R^{\pm}\), \(B\in GL_4(K)\).
For \(A_r\) one can take \(R^{\pm}=\{\pm P_\sigma:
\sigma\in S_{r+1}\}\); for the exceptional real models
the exact root normalizer is the one retained in HD-I23.

Keep the whole finite preimage
\[
 \widehat G=\{(A,B)\in R^{\pm}\times GL_4(K):
                         A\otimes B\in CGC^{-1}\}.
                                                               \tag{R19.2}
\]
Its projection to the original \(G\) is surjective.
An identity tensor has scalar factors, and the scalar
matrices in the real finite orthogonal \(R^{\pm}\)
are exactly \(\pm I\). Its kernel is therefore
\(\{(I,I),(-I,-I)\}\), and
\(|\widehat G|=2|G|\). In particular every factor
matrix has finite order; no arbitrary infinite scalar
extension has been made.

Choose original sphere inertia generators \(g_1,\ldots,g_b\),
whose product is one, and choose any lifts
\((A_j,B_j)\in\widehat G\). Their products satisfy
\[
 \prod_jA_j=\varepsilon I_r,\qquad
 \prod_jB_j=\varepsilon I_4,\qquad\varepsilon\in\{1,-1\}.
                                                               \tag{R19.3}
\]
The two groups generated by their factor matrices are
irreducible. For example a proper invariant subspace of
\(M\) would give the proper invariant subspace
\(C^{-1}(W\otimes M')\) for every original \(g_j\),
and hence for \(G\). The proof on \(W\) is the same.
Adding the scalar kernel in (R19.2) does not change
these invariant subspaces. Write
\(c_j=\operatorname{rank}(A_j\otimes B_j-I)\). The
original Euler equation and period dimension give
\[
                  \sum_jc_j=8r+6,\qquad
                          b\ge12,\quad c_j\ge4.         \tag{R19.4}
\]

\paragraph{2. A scalar-product tuple and its image subspaces.}
We first prove the exact linear-algebra bound used below.
For any irreducible tuple of semisimple operators
\(B_1,\ldots,B_b\) on a complex space of dimension
\(d\ge2\), with scalar product, put
\(\delta_j=d-\max_\beta\dim\ker(B_j-\beta I)\).
Then
\[
                             \sum_j\delta_j\ge d+1.     \tag{R19.5}
\]
To prove it, omit the scalar entries and rescale each
remaining \(B_j\) by an eigenvalue of largest
multiplicity, giving \(C_j\) with
\(\operatorname{rank}(C_j-I)=\delta_j\).
Their product is still a scalar \(\lambda I\).
Put \(U_j=\operatorname{im}(C_j-I)\).
The sum \(U=\sum U_j\) is invariant under every
\(C_i\), because \(C_i u=u+(C_i-I)u\in U+U_i\).
It is nonzero, so irreducibility makes \(U=V\), and
\(\sum\delta_j\ge d\).

Equality would give the actual direct sum
\(V=\bigoplus_{j=1}^k U_j\). In its block coordinates,
\(C_j-I\) has nonzero entries only in block row \(j\).
The last block row of \(C_1\cdots C_k\) equals the
last block row of \(C_k\); the product \(\lambda I\)
therefore makes \(C_k\) diagonal by blocks, equal to
\(\lambda I\) on \(U_k\) and the identity on the
other summands. Remove it from the product. The same
argument with row \(k-1\), and then successively with
the previous rows, makes every \(C_j\) the corresponding
diagonal block operator. Thus every \(U_j\) is invariant
under the tuple. Each \(\delta_j\) lies in
\(\{1,\ldots,d-1\}\), so equality has at least two
nonzero blocks, contradicting irreducibility. This
proves (R19.5), including the scalar-product case with
all scalar phases retained.

For our four-dimensional factor, let \(k\) be the number
of projectively nontrivial \(B_j\), and put
\(S=\sum\delta_j\). Then \(S\ge5\) and \(k\le S\).
For a fixed eigenvalue of \(A_j\), its contribution to
\(\ker(A_j\otimes B_j-I)\) is at most
\(4-\delta_j\) times that eigenvalue multiplicity.
Summing over the actual eigenspaces of \(A_j\) gives
\(c_j\ge r\delta_j\). The other \(b-k\) original
inertias still have \(c_j\ge4\). Therefore
\[
 8r+6=\sum_jc_j\ge rS+4(b-k)
                         \ge(r-4)S+48.                 \tag{R19.6}
\]
For \(r=6\) the right side is at least58, exceeding54.
For \(r=7\) it is at least63, exceeding62.
These two ranks are excluded without using the real
structure of the first factor.

\paragraph{3. The exact rank-eight equality range.}
Now let \(r=8\). Formula (R19.6) gives
\(70\ge4S+48\), hence \(S=5\). Also
\(70\ge40+4(b-k)\), with \(b\ge12\), \(k\le5\).
It follows that \(k=5\) and \(b=12\): either \(k\le4\)
or \(b\ge13\) would give at least72. All five
projectively nontrivial entries have \(\delta_j=1\).
Their lower bounds contribute40 and the seven scalar
entries contribute28. The total excess over these
bounds is exactly two.

Each of the five nonscalar \(B_j\) has eigenvalues
\(\beta_j\) of multiplicity three and
\(\alpha_j\ne\beta_j\) of multiplicity one. The
matrices have finite order, so these are unit roots.
If \(\beta_j\) were nonreal, the real orthogonal
\(A_j\) would have multiplicity
\(a=\dim\ker(A_j-\beta_j^{-1}I)\le4\), since its
conjugate eigenvalue has the same multiplicity. The
other matching multiplicity is at most \(8-a\).
Thus the tensor fixed dimension would be at most
\(3a+(8-a)\le16\), and \(c_j\ge16\).
Its excess over the bound8 would already exceed two.
Consequently \(\beta_j=\pm1\).

For a scalar entry \(B_j=\beta_jI_4\), a nonreal
\(\beta_j\) would likewise give tensor fixed dimension
at most16, hence \(c_j\ge16\), exceeding its bound4
by more than two. Its scalar is therefore also \(\pm1\).
Absorb every such sign into the real factor:
\[
 A'_j=\beta_j A_j,\qquad B'_j=\beta_j^{-1}B_j,
                   \qquad A'_j\otimes B'_j=A_j\otimes B_j.
                                                               \tag{R19.7}
\]
These are actual matrices in the same finite preimage;
their products are still scalar signs. The scalar
entries now have \(B'_j=I_4\), while the other five
have spectrum \((1,1,1,\lambda_j)\), \(\lambda_j\ne1\).
The generated groups on both factors remain irreducible,
since multiplication by scalar signs preserves exactly
the same invariant subspaces.

For each scalar entry put \(q_j=\operatorname{rank}(A'_j-I)\).
Its original inertia is nonidentity, so \(q_j\ge1\)
and \(c_j=4q_j\). The excess bound two forces
\(q_j=1\). Each such \(A'_j\) is a real orthogonal
reflection: a finite-order real orthogonal matrix with
fixed hyperplane has remaining eigenvalue \(-1\).

For each of the other five entries, write
\(q_j=8-\dim\ker(A'_j-I)\) and
\(t_j=\dim\ker(A'_j-\lambda_j^{-1}I)\le q_j\).
The exact tensor eigenspace formula is
\[
                  c_j=8+3q_j-t_j\ge8+2q_j.              \tag{R19.8}
\]
An excess at most two allows \(q_j=0\), giving
\(A'_j=I\) and \(c_j=8\), or \(q_j=1\).
In the latter case \(A'_j\) is a real reflection;
an excess at most two then forces \(t_j=1\), hence
\(\lambda_j=-1\), and \(c_j=10\). All other
possibilities exceed the allowed excess. Since the
total in (R19.4) is70 rather than68, exactly one of
these five entries is the latter case and the other
four have \(A'_j=I\).

The actual \(A'_j\)-tuple therefore consists of eight
real reflections and four identities, has scalar
product, and generates an irreducible group on the
eight-dimensional \(W\). Apply (R19.5) to this very
tuple with \(d=8\). Each reflection has largest
eigenspace dimension seven, so its \(\delta\) is one;
each identity has \(\delta=0\). Their sum is eight,
whereas (R19.5) requires at least nine. This is a
contradiction and excludes \(r=8\).

\paragraph{4. Retained source and remaining multiplicity-four case.}
All inequalities above are on the original inertia
matrices transported by (R19.1). Their product relation,
Euler characteristic, finite cover and primitive
Jacobian corner have never been replaced. In particular
the possible diagonal sign kernel in (R19.2) has been
retained in every tuple relation. The original positive
coefficient form is transported by \(C\); comparison
with a prescribed invariant form changes it by a
positive rational scalar, which cancels in the normalized
matrix coefficient exactly as in HD-I23 and HD-R18.
Thus the contradiction applies to the prescribed
polarized sixfold factor, not merely to a complex
representation with a matching character.

Combined with the supplied imprimitive reflection
case and the complete reflection-constituent list
in HD-I23, this removes all constituent ranks below
nine from the multiplicity-four alternative. Its
remaining actual form is the rational augmentation
factor of rank \(r\ge9\), tensored with the actual
four-dimensional factor; the original overgroup and
its sphere tuple remain to be analyzed there.
